\section{Related Work}\label{sec:related}

\paragraph{Exploration and exploitation in groups.}
\citet{march1991exploration} modelled an organisation whose members learn from a shared code and showed that
faster mutual learning raises average knowledge early and lowers it in the long run, because the code absorbs
the members' beliefs before the best belief has been found. \citet{lazer2007network} placed agents on an NK
landscape \citep{levinthal1997adaptation} and varied how densely they were connected: dense networks spread the
current best solution quickly, which raises short-run performance and lowers the long-run outcome, while
sparse networks keep diversity alive long enough to find better regions. \citet{fang2010balancing} found that
semi-isolated subgroups that occasionally exchange information outperform both fully connected and fully
isolated groups. Experiments with people reproduce the pattern. \citet{mason2012collaborative} let networked
participants search a landscape with a hidden peak and found that efficient networks did better on average
but could converge prematurely, and \citet{derex2016partial} found that partially connected groups accumulated
more cultural innovation than fully connected ones. Our periodic-board condition and our private-rounds
ladder are direct analogues of the sparse and semi-isolated structures in this literature, run with LLM agents
instead of simulated or human ones.

\paragraph{Search on rugged landscapes with people.}
\citet{ederer2013pay} designed the lemonade-stand game to test whether paying for early performance deters
exploration of a strategy that is better in the long run but worse at first. Participants who were paid every
period stayed near the recommended setup; participants whose early failures were tolerated explored more and
more often found the better location. \citet{billinger2014search} ran a related experiment on NK landscapes
and found that people search less than the landscape warrants, especially after early success. We keep the
game, the valley and the letter, and change the question from how the individual is paid to what the
individual can see of others.

\paragraph{LLM agents and exploration.}
Frozen LLMs explore poorly in bandit problems: they fail to return to the best arm, or play nearly uniformly,
except under engineered prompts \citep{krishnamurthy2024explore}, and their failures are greedy
\citep{schmied2025greedy}. Their exploration is also shaped by what the arms are called:
\citet{austin2026semantic} vary the labels of a bandit and find that semantic priors bias which arms are
tried. Our alien framing follows that idea. It keeps every number in the world and changes only the words,
so that any difference between framings is attributable to the labels and to the recommended starting setup,
not to the landscape.

\paragraph{Communication between LLM agents.}
Generative-agent systems let LLM agents observe and talk to one another \citep{park2023generative}, and
\citet{park2026scaling} study the minimal case we start from: identical agents, one shared workspace, no roles
and no orchestrator. They find that sharing beats the best of $k$ independent agents when a verifier ranks
partial progress correctly and that the advantage disappears when it cannot. Our world is the case they leave
open: a verifier that is correct about the current round and misleading about the direction of the best one.
